% !TeX root = main_aip.tex
\newpage
\section{Dynamics}

\subsection{Description of Dynamics}

The dynamics of FPS:

\begin{enumerate}
\item It is purely geometric 
	\begin{itemize}
	\item that all calculations are founded on the area of overlap of particles.
	\end{itemize}
\item Particle collisions are interpenetrating
	\begin{itemize}
	\item Particles stay in contact during compression and rebound.
	\item Allows for time to perform in contact operations
	\item Creates the need for storing momentum as a replacement for internal energy.
	\end{itemize}
\item It is instantaneous
	\begin{itemize}
	\item  meaning that all force and momentum calculations occur at a single time step. 
	\item without regards to previous values except prior position and velocity.
	\end{itemize}
\item It is energyless
	\begin{itemize}
	\item energy is a diagnostic value only
	\end{itemize}
\item All particle data is source owned.
	\begin{itemize}
	\item creates a source owned system that is parallel friendly and fits with gPCD.
	\end{itemize}
\end{enumerate} 

\subsection{Geometry-to-Motion Causality}

The Field Particle System (FPS) may be viewed as a
geometry-to-motion system. The purpose of the methodology
is not to solve governing equations directly, but rather
to transform geometric relationships into particle motion
through a sequence of local interactions. At every
timestep, the current geometric configuration of the
system determines the interaction state of each particle.
The resulting interactions generate forces, forces produce
accelerations, accelerations modify motion, and the
resulting motion creates a new geometric configuration.
The process then repeats.

The causal chain begins with particle geometry. Particle
position, size, shape, orientation, and boundary
relationships collectively define the geometric state of
the system. This state determines whether particles are
separated, approaching contact, overlapping, bonded, or
interacting with boundaries.

When particles enter an overlapping configuration, an
interaction state is established. The overlap region is a
purely geometric quantity that measures the extent of
particle interpenetration and defines a temporary domain
within which particle interactions may occur. Depending
upon the simulation model, this interaction state may
support force generation, momentum transfer, bonding,
species exchange, thermal exchange, or other particle-
level processes.

The overlap geometry is subsequently converted into a
contact area. The contact area serves as a geometric
measure of interaction strength and provides the basis for
force generation. Larger overlap regions generally produce
larger contact areas and therefore stronger interactions.

The force-generation model converts contact area into
force magnitude through the particle-pair stiffness
relationship. The resulting force acts upon the source
particle and produces an acceleration according to the
particle mass.

Acceleration modifies particle velocity through numerical
integration. The updated velocity then modifies particle
position, producing a new geometric configuration for the
next timestep.

The complete causal chain may therefore be summarized as

\begin{verbatim}
Geometry
    ↓
Interaction State
    ↓
Overlap
    ↓
Contact Area
    ↓
Force
    ↓
Acceleration
    ↓
Velocity
    ↓
Position
    ↓
New Geometry
\end{verbatim}

An important consequence of this formulation is that
geometry acts as the fundamental source of causality
within FPS. Motion changes geometry, geometry creates
interactions, and interactions generate new motion.
Consequently, the methodology forms a closed causal loop
in which the evolution of the system is governed by the
continuous conversion of geometry into motion and motion
into geometry.

This viewpoint distinguishes FPS from methodologies in
which energy, pressure, constitutive laws, or other
macroscopic quantities serve as the primary drivers of
system evolution. Within FPS, such quantities may be
observed, measured, or derived; however, the underlying
causal mechanism remains geometric interaction.

\subsubsection{Geometry-to-Motion Causality Plot}
\begin{verbatim}


\end{verbatim}

\subsection{Purely Geometric Interactions}

The Field Particle System (FPS) employs a purely geometric
interaction formulation in which particle behavior is
determined exclusively from geometric relationships between
particles and boundaries. Interaction quantities are derived
from geometric measures such as overlap, contact orientation,
separation distance, particle size, and boundary geometry.

The term \emph{purely geometric} does not imply the absence of
forces. Rather, it indicates that all forces within the
framework originate from geometric relationships. Force
generation is therefore a consequence of particle geometry and
spatial configuration rather than the direct application of
phenomenon-specific governing models.

Unlike many traditional simulation methodologies, FPS does not
require equations of state, constitutive relationships,
turbulence models, shock-capturing algorithms, transition
criteria, or other phenomenon-specific formulations to govern
particle interactions. Instead, a common set of geometric
interaction rules is applied throughout the simulation domain
regardless of the macroscopic behavior that may emerge.

As a consequence, the same interaction methodology governs all
particle species and simulation scenarios. Complex system-level
behavior arises from the collective evolution of local
geometric interactions rather than from specialized models
introduced to reproduce particular physical phenomena.

The purely geometric formulation is a fundamental characteristic
of FPS and forms the basis for the framework's ability to apply
a single interaction methodology across a diverse range of
simulation environments.
```


\subsection{Instantaneous Dynamics}

The Field Particle System (FPS) employs an instantaneous
dynamics formulation in which particle response is determined
from the present geometric state of the system. Interactions do
not require a stored history of previous contacts, accumulated
collision states, or persistent memory of prior particle
interactions.

At each simulation step, forces and accelerations are evaluated
from the current geometric configuration of the particles. The
resulting dynamic response is applied during the same timestep,
so the particle motion is governed by the geometry that exists
at the moment of evaluation.

The only exception occurs when a particle is initialized in an
overlapping or colliding state. In that case, the initial contact
condition may require stored information to identify the particle
as already participating in a collision at the start of the
simulation. Apart from this initialization condition, the FPS
dynamics do not require memory of previous contacts in order to
compute the current response.

The term \emph{instantaneous} does not imply infinite propagation
speed or instantaneous communication across the simulation
domain. Rather, it indicates that the local response of each
particle is computed directly from the current particle geometry
and applied immediately within the timestep.

\subsection{Energyless Dynamics}

The Field Particle System (FPS) employs an energyless
dynamics formulation in which energy is not treated as a
governing quantity of the simulation. Particle motion is
determined directly from geometric interactions and the
resulting forces without requiring the storage, transport,
exchange, minimization, or conservation of energy.

Within FPS, energy is used only as a diagnostic quantity
during analysis. It may be computed from the evolving
particle configuration to evaluate or interpret system
behavior, but it does not participate in the determination
of particle motion.

This distinction is also reflected in the executable
implementation. When the method is exported to GLSL for
GPU execution, diagnostic quantities are removed from the
runtime code. The exported shader implementation therefore
evolves particle states without computing or storing
energy-related diagnostic terms.

The framework does not impose energy conservation as a
governing constraint. Momentum may exhibit conservation
properties arising from the symmetry of the interaction
rules, whereas energy is not required to remain constant.
As a result, energy may be created or destroyed during
system evolution while still remaining useful as an
offline diagnostic measure.

Accordingly, FPS treats energy as an analytical observable
rather than a governing variable. The simulation is driven
by geometric interactions and their associated forces,
while energy diagnostics exist only outside the core
runtime dynamics.

```latex
\subsection{Force Generation}

Within FPS, forces arise directly from the geometric
relationship between interacting particles. When two
particles overlap, the intersection region defines a
contact area that characterizes the strength of the
interaction. Larger overlaps produce larger contact
areas and therefore stronger interactions, while
particles that are not in contact generate no force.

For a particle pair $(i,j)$, the overlap geometry is
used to compute the contact area $A_{ij}$. The force
magnitude is then obtained from

\begin{equation}
F_{ij}=q_{ij}A_{ij},
\end{equation}

where $q_{ij}$ is the pair stiffness coefficient and
$A_{ij}$ is the overlap area. The pair stiffness
controls the interaction strength while the overlap
area determines how strongly the particles interact
within the current geometric configuration.

The resulting force acts along the line connecting the
particle centers. Let

\begin{equation}
\mathbf{n}_{ij}
=
\frac{\mathbf{x}_{j}-\mathbf{x}_{i}}
{\left\|\mathbf{x}_{j}-\mathbf{x}_{i}\right\|}
\end{equation}

denote the unit vector directed from particle $i$ to
particle $j$. The force applied to particle $i$ is

\begin{equation}
\mathbf{F}_{ij}
=
F_{ij}\mathbf{n}_{ij}.
\end{equation}

Because the force is derived entirely from overlap
geometry and particle configuration, no equations of
state, constitutive relationships, turbulence models,
shock-capturing algorithms, or other phenomenon-
specific governing equations are required to generate
particle interactions.

The total force acting on a particle is obtained by
summing the contributions from all interacting
neighbors. These forces are subsequently converted to
accelerations and used by the integration procedure to
advance the system through time.

```latex id="2xv7fr"
\subsection{Numerical Integration}

The FPS interaction model is independent of the numerical
integration procedure used to advance particle states.
Once forces have been computed from the geometric
configuration of the system, any suitable integration
scheme may be employed to update particle positions and
velocities.

The current implementation utilizes semi-implicit Euler
integration because of its simplicity, computational
efficiency, and suitability for massively parallel GPU
execution. For each particle, acceleration is first
computed from the total applied force

\begin{equation}
\mathbf{a}_i=\frac{\mathbf{F}_i}{m_i}.
\end{equation}

Velocity is then updated according to

\begin{equation}
\mathbf{v}_i^{\,n+1}
=
\mathbf{v}_i^{\,n}
+
\mathbf{a}_i^{\,n}\Delta t,
\end{equation}

followed by the position update

\begin{equation}
\mathbf{x}_i^{\,n+1}
=
\mathbf{x}_i^{\,n}
+
\mathbf{v}_i^{\,n+1}\Delta t.
\end{equation}

Because the updated velocity is used during the position
update, the method provides improved stability compared
to explicit Euler integration while retaining a simple
computational structure.

It should be noted that the FPS methodology does not
require semi-implicit Euler integration specifically.
Alternative integration procedures may be employed
without modifying the underlying interaction model,
provided that particle positions and velocities are
advanced consistently with the computed forces.

\subsubsection{Drift and Excursion}

When evaluating the long-term behavior of a numerical
integration scheme, it is useful to distinguish between
drift and excursion.

Drift refers to the systematic accumulation of numerical
error that produces a persistent trend in a measured
quantity. Excursion refers to a temporary deviation from
an expected value that remains bounded and subsequently
reverses direction.

Evaluation of the FPS implementation indicated that
measured deviations were dominated by bounded excursions
rather than persistent drift. Quantities exhibiting
temporary departures from their expected values were
observed to oscillate about a mean state without
demonstrating continuous growth over time.

This behavior is consistent with the use of a symplectic
integration scheme. Symplectic methods do not generally
eliminate numerical error; however, they tend to prevent
the systematic accumulation of error associated with
long-term drift. As a consequence, errors often manifest
as bounded excursions rather than continuously growing
deviations.

For the simulations examined in this work, the observed
behavior suggests that numerical error remains bounded
over extended simulation periods, thereby reducing the
likelihood of unphysical long-term drift.


```latex
\subsection{Interpenetrating Contacts}

The Field Particle System (FPS) employs an interpenetrating
contact model in which particles are permitted to occupy a
common geometric region during interaction. Unlike rigid-body
collision formulations that attempt to prevent overlap through
instantaneous collision resolution, FPS allows overlap to
develop naturally and uses the resulting overlap geometry as
the basis for force generation.

Within the framework, overlap is not interpreted as a
simulation error requiring immediate correction. Instead,
overlap is regarded as a physically meaningful geometric
state that quantifies the strength of interaction between
particles. As particles move into one another, the overlap
region grows, increasing the geometric contact area and
therefore increasing the force generated by the interaction.
As particles separate, the overlap region decreases and the
resulting force diminishes.

Interpenetration is not merely a mechanism for force generation. 
It establishes a finite-duration interaction state during which particles 
may exchange information, transfer momentum, form bonds, or participate in 
other interaction processes. The overlap region therefore serves as both a 
geometric measure of contact and a temporary domain in which 
particle-level interactions may occur.

This interpretation differs fundamentally from many
traditional collision models. In rigid collision systems,
the primary objective is often to eliminate penetration and
restore a non-overlapping configuration. FPS adopts the
opposite viewpoint. Penetration is not removed directly.
Rather, penetration generates forces that influence future
particle motion. The reduction of overlap therefore occurs
through dynamic evolution of the system rather than through
explicit geometric correction.

Interpenetration also provides a continuous description of
contact. Traditional collision systems frequently treat
contact as a discrete event occurring at a particular
instant in time. FPS instead treats contact as a finite
process extending over multiple timesteps. A particle pair
may enter contact, increase overlap, reach a maximum
interaction state, and subsequently separate, all while
remaining governed by the same interaction rules.

The overlap region therefore serves as an intermediate
quantity connecting geometry and dynamics. Particle
configuration determines overlap. Overlap determines
contact area. Contact area determines force magnitude.
Force determines acceleration and motion. Through this
sequence, geometric interpenetration becomes the mechanism
by which particles influence one another.

An important consequence of this formulation is that the
same interaction mechanism applies throughout the simulation
domain. No distinction is required between collision,
contact, compression, rebound, or separation phases.
Particles interact through a common geometric process, and
the resulting behavior emerges from the evolution of the
overlap geometry.

Within FPS, interpenetration is therefore not a numerical
artifact but a fundamental component of the dynamics. The
ability of particles to occupy a shared geometric region
provides the foundation for force generation and enables
the methodology to convert geometry directly into motion.
```
\subsection{Source-Owned Dynamics}

The Field Particle System (FPS) employs a source-owned
dynamics formulation in which each particle is solely
responsible for the computation and modification of its
own state. Although a particle may observe neighboring
particles and use their properties during force
evaluation, it never directly modifies the state of
another particle.

Within the methodology, the particle being updated is
referred to as the source particle. Neighboring particles
participating in the interaction are referred to as
target particles. During force evaluation, a source
particle may read information from any number of target
particles; however, all resulting state changes remain
local to the source particle.

Consequently, interactions are not treated as pair-owned
operations. The interaction between two particles is
evaluated independently from the perspective of each
participant. Particle \(i\) computes its response to
particle \(j\), and particle \(j\) computes its response
to particle \(i\) during its own update procedure.
Neither particle writes to the state of the other.

This ownership model eliminates the need for pairwise
state updates, shared modification of particle data, and
many forms of synchronization commonly encountered in
parallel particle systems. Each particle remains the sole
owner of its position, velocity, momentum, bonding state,
and other dynamic quantities throughout the simulation.

Source ownership also provides a natural framework for
future interaction processes. Momentum transfer, bonding,
species exchange, thermal exchange, and other particle
interactions may be evaluated from the perspective of the
source particle while preserving exclusive ownership of
particle state.

Within FPS, source ownership is therefore not merely a
parallel implementation strategy. It is a fundamental
property of the dynamics that defines how information
flows through the system and how particles interact with
one another.
